Two principles govern this section, and between them they settle most of it.

\begin{quote}
\textbf{The record is the row.} A module's record type has exactly the columns
of its table --- no hidden columns and no projections --- so that
$\llbracket \texttt{M.t} \rrbracket$ \emph{is} the row type of $M$'s table.
Accessors are derived navigation on top of it.

\textbf{The store is explicit.} Every function that reads rows takes the
database as its first argument, so every generated function is a function of
its arguments.
\end{quote}

\noindent
The generated interface therefore opens with an abstract handle:

\begin{lstlisting}[style=ml]
type db
\end{lstlisting}

\noindent
How a \texttt{db} is obtained is a question about storage rather than about
lowering, and is not specified here.

\subsection*{S21.1. Abstractness of \texttt{id}}

How much of the identifier type the generated signature exposes.

\options
\begin{enumerate}[label=\textbf{(\alph*)}]
  \item \textbf{Fully abstract} --- \texttt{type id}, with no operations. The
        only thing a consumer can do with an identifier is pass it back to
        \texttt{get}. \selected
  \item \textbf{Abstract with an operations bundle} --- \texttt{equal},
        \texttt{compare}, \texttt{to\_string}, \texttt{hash}.
  \item \textbf{\texttt{type id = private int}}, which coerces freely to
        \texttt{int} but cannot be forged.
  \item \textbf{\texttt{type id = int}}, transparent.
\end{enumerate}

\decision{Fully abstract, as written in the design note.}

\begin{lstlisting}[style=ml]
module Root : sig
  type id
  type t = { id : id; a : int }
  val get : db -> id -> t
end
\end{lstlisting}

\rationale{This makes \texttt{Root.id} and \texttt{B.id} distinct types, so an
identifier belonging to one table cannot be passed where another's is
expected --- an easy mistake to make when every identifier is an integer
underneath. It is also faithful to the design note's framing of these columns
as generated rather than part of the document.}

\limitation{A consumer cannot print an identifier, compare two, hash one, or
use one as a map key. Both (b) and (c) remain available, and adopting either
later is backward compatible: (b) only adds functions, and (c) only weakens an
abstract type to a private abbreviation.}

\subsection*{S21.2. The store behind \texttt{get}}

\options
\begin{enumerate}[label=\textbf{(\alph*)}]
  \item \textbf{An implicit ambient store}, as the design note writes it.
  \item \textbf{An explicit handle} --- \texttt{val get : db -> id -> t}.
        \selected
  \item \textbf{A functor over a storage signature}.
  \item \textbf{The store is the effect handler}, supplying rows in response
        to a performed effect.
\end{enumerate}

\decision{An explicit handle. Every function that reads rows takes a
\texttt{db} as its first argument.}

\rationale{Under an ambient store no generated function is a function of its
arguments: \texttt{get : id -> t} reads hidden mutable state, two calls need
not agree, and the signature does not admit that it depends on anything but
its identifier. Threading the handle makes each generated function pure and
its denotation immediate --- $\llbracket \texttt{get} \rrbracket\ d\ i$ is the
row named by $i$ in $d$ --- which is what any equivalence statement about
these modules has to quantify over.

This diverges from the design note, which writes the ambient form. The
divergence is deliberate and is the price of a specification that can be
stated without an implicit parameter running through every clause.}

\paragraph{\texttt{get} is total.}
It requires no \texttt{option} and cannot fail. Identifiers are abstract
(S21.1), so the only way to obtain one is from a module that already holds it;
and foreign-key constraints are materialized, so every identifier in
circulation names a row that exists. Totality follows from those two decisions
together rather than being assumed.

\limitation{Every accessor gains a parameter, which makes navigation chains
noisier to write: \texttt{B.c d (Root.b d r)} rather than
\texttt{B.c (Root.b r)}. Option (d) remains the eventual destination if
Phase~2's handlers come to supply rows, at which point the handle could be
dropped again.}

\subsection*{S21.3. The presence marker in the record type}

S13 records presence in the parent when an array-valued path is partial, so
that the accessor can return an \texttt{option}. This asks whether that marker
is visible in the record.

\options
\begin{enumerate}[label=\textbf{(\alph*)}]
  \item \textbf{Only through the accessor.} The record holds nothing.
  \item \textbf{In the record as well}, as a field mirroring the column.
        \selected
  \item \textbf{In the record only}, with the accessor returning a plain list.
\end{enumerate}

\decision{The marker appears in the record, because it is a column and the
record is the row.}

\begin{lstlisting}[style=ml]
module Root : sig
  type id
  type t = { id : id; items_present : bool }
  val get : db -> id -> t
  val items : db -> t -> Item.t list option
end
\end{lstlisting}

\rationale{Follows directly from the governing principle. Keeping the record
in exact correspondence with the table means the generated type needs no
separate description: it is the row, and nothing has to be said about which
columns were withheld.}

\limitation{A consumer now has two representations of one fact --- the
\texttt{items\_present} field and the accessor's \texttt{option} --- which can
be read inconsistently, and \texttt{items\_present : bool} beside an
\texttt{option}-returning accessor is the least attractive consequence of the
principle. Option (c) is excluded regardless: it withholds the \texttt{option}
and moves the obligation to check onto the caller.}

\subsection*{S21.4. Foreign key field naming}

A foreign key appears twice: as the record field holding the child's
identifier, and as the accessor that follows it. Naming both after the JSON
member makes \texttt{r.b} and \texttt{b d r} look nearly identical while having
different types.

\options
\begin{enumerate}[label=\textbf{(\alph*)}]
  \item \textbf{Both named \texttt{b}}, as the design note's code writes it.
  \item \textbf{Field \texttt{b\_id}, accessor \texttt{b}} --- matching the
        relational column. \selected
  \item \textbf{Field \texttt{b}, accessor renamed} to \texttt{get\_b}.
  \item \textbf{Hide the foreign key from the record.}
\end{enumerate}

\decision{The field is named for the column, \texttt{b\_id}; the accessor
keeps the member's name.}

\begin{lstlisting}[style=ml]
module Root : sig
  type id
  type t = { id : id; a : int; b_id : B.id }
  val get : db -> id -> t
  val b : db -> t -> B.t          (* follows b_id into table b *)
end
\end{lstlisting}

\rationale{This follows from the governing principle rather than being an
independent choice: the column is named \texttt{b\_id}, so the field is too.
It removes the clash for free.

The design note is not internally consistent here --- its relational table
names the column \texttt{b\_id} while its generated record names the field
\texttt{b} --- so this decision follows the note's table where S21.1 followed
its code. Option (d) is excluded by the principle, since a foreign key is a
column.}

\subsection*{S21.5. Entry point into the modules}

Every accessor starts from a value of type \texttt{t}, and \texttt{get}
requires an identifier obtainable only from a module already holding one.
Nothing in the design note produces the first row.

\options
\begin{enumerate}[label=\textbf{(\alph*)}]
  \item \textbf{\texttt{val all : db -> t list}} on every generated module.
        \selected
  \item \textbf{\texttt{val all : db -> t Seq.t}} --- lazy, so the table is
        not materialized.
  \item \textbf{An entry point on root tables only.}
  \item \textbf{A separate load step} returning a handle carrying the roots.
\end{enumerate}

\decision{Every generated module exposes \texttt{val all : db -> t list},
returning every row of its table.}

\rationale{The simplest and most uniform of the four: every table is reachable
without navigating to it. Option (c) would leave whole tables reachable only
by navigation, ruling out queries that range over a table independently of its
parents. Option (d) is subsumed by S21.2, which already supplies the handle.}

\limitation{\texttt{all} materializes the entire table, which is what the
design note's \S3.2 argues against. Option (b) is the natural revision, though
changing \texttt{t list} to \texttt{t Seq.t} would break consumers rather than
extend them.}

\subsection*{S21.6. The child's back-reference in the record type}

A child table from S13 carries \texttt{parent\_id} and \texttt{idx} columns.
This asks whether they appear in the child's record.

\options
\begin{enumerate}[label=\textbf{(\alph*)}]
  \item \textbf{Hide both}, as encoding artifacts.
  \item \textbf{Expose both as fields.} \selected
  \item \textbf{Expose via accessors} --- \texttt{val parent : db -> t ->
        Root.t} --- making navigation symmetric.
  \item \textbf{Expose \texttt{idx} only.}
\end{enumerate}

\decision{Both appear in the record, because both are columns.}

\begin{lstlisting}[style=ml]
module Item : sig
  type id
  type t = { id : id; root_id : Root.id; idx : int; sku : string; qty : int }
  val get : db -> id -> t
  val all : db -> t list
end
\end{lstlisting}

\rationale{The governing principle again. It also makes S21.5's choice pay
off: with \texttt{all} returning every item from every parent flattened
together, a consumer could otherwise not tell which parent a row belonged to
or where it sat in the array.

Note that \texttt{idx} is not purely an artifact in the way a generated
identifier is. Array order is genuine document data, and \texttt{idx} is how
it is represented, so hiding it would discard something the document
stated.}

\limitation{Generated identifiers, back-references and ordering indices are
now all public. The interface therefore exposes the encoding, and a consumer
can write code that depends on it --- for instance sorting by \texttt{idx}
directly rather than relying on the order \texttt{items} returns.}
